% ============================================================
\section{Background}
% ============================================================

\subsection{The longitudinal measurement gap in spine biomechanics}

The aetiology of non-specific \ac{LBP} --- the leading cause of years lived
with disability globally \cite{hartvigsen2018} --- is rarely a single traumatic event.
It reflects the cumulative effect of repetitive microtrauma: sustained flexed postures at
a desk, asymmetric load bearing in care work, or hours spent in suboptimal sleep
positions \cite{burdorf1997}.
Identifying the specific movement habits that drive individual pathology requires
measurement instruments capable of operating continuously in real-world environments,
not just for the duration of a laboratory session.

Contemporary assessment options divide into four categories, each with a decisive
limitation for this purpose.

\subsection{The clinical diagnostic blind spot}

Radiological imaging --- conventional X-ray, \ac{MRI}, and \ac{CT} --- provides
indispensable high-resolution information about bony structures and soft tissue.
From a biomechanical perspective, however, these modalities yield only a static
snapshot.
A patient lying in an MRI scanner is in a supine, fully unloaded position in which
gravitational forces on the intervertebral discs are minimal.
The resulting images show the morphological state --- a compressed disc prolapse,
a degenerative spinal canal stenosis, or facet-joint arthrosis --- but contain no
causal information about the dynamic mechanical exposure in daily life that drove
that tissue degeneration \cite{natarajan2023}.

Non-invasive surface measurement systems such as rastersterography and mechanical
measuring wheels (e.g., the Idiag M360 / SpinalMouse) capture sagittal and frontal
curvature with high precision and without radiation \cite{spinalMouse}.
These instruments are well suited to supervised in-clinic assessments and therapeutic
follow-up, but by design they yield a single static snapshot taken during a brief,
controlled examination session.

Once a patient leaves the clinic and returns to daily life, a large diagnostic blind
spot opens.
The exposure events most relevant to pathogenesis --- hours of sustained lumbar flexion
at a PC workstation, asymmetric lifting of patients from low hospital beds, monotonous
load carrying in logistics work, or prolonged adverse sleep positions --- are entirely
invisible to objective measurement.
In this vacuum, research frequently relies on Patient-Reported Outcome Measures
(\acp{PROM}).
These instruments carry well-documented susceptibility to recall bias, social
desirability effects, and cognitive distortion, limiting their validity for
establishing causal relationships between daily postural habits and pain
episodes \cite{natarajan2023}.

\subsection{Optical motion-capture systems}

Optoelectronic systems such as Vicon, Qualisys, and OptiTrack track retro-reflective
markers with sub-millimetre precision --- Ehara et al.\ report RMSE values around
0.154\,mm for Vicon \cite{vicon_rmse} --- and have served as the biomechanical reference
standard for decades.
Three structural constraints prevent deployment outside a dedicated capture volume.

First, the technology requires a calibrated array of synchronised infrared cameras;
transporting and recalibrating this array for a building site, a care ward, or a
patient's bedroom is impractical.
Second, any occlusion of the marker-to-camera line of sight --- by furniture, protective
equipment, a patient in a hospital bed, scaffolding, or ordinary clothing --- causes
immediate data loss or trajectory interpolation artefacts.
Third, the requirement to affix dozens of spherical markers to anatomical landmarks on
the skin (or tight-fitting specialised clothing) places participants in a manifestly
artificial laboratory environment.
This \emph{lab effect} compromises ecological validity: participants modify their
natural movement behaviour under camera observation and equipment constraints, so the
recorded postures no longer represent their true everyday kinematics.

Marker-less deep-learning pose estimation reduces setup burden but retains
occlusion sensitivity, is vulnerable to changes in viewpoint and illumination,
and cannot resolve inter-vertebral angles with the precision required for clinical
spine metrics.

\subsection{Standard inertial measurement units}

A single \acs{IMU} placed on the thorax measures the global inclination of the upper torso
relative to gravity.
This tells a researcher how far the torso is tilted, but not \emph{where along the
	spine} the flexion occurs --- whether the pelvis has rotated on an extended lumbar spine
or whether a pathological thoracic kyphosis is responsible.
This distinction is clinically essential: the biomechanical consequences of hip-hinge
flexion with a neutral lumbar spine are entirely different from those of lumbar and
thoracic hyperkyphosis under load \cite{burdorf1997}.

Arrays of multiple \acp{IMU} can partially resolve segmentation, but the approach
scales poorly.
Full-body suits carrying up to 17 individual Bluetooth nodes are thermally
uncomfortable for participants, prone to displacement during occupational or athletic
activity, and difficult to synchronise reliably over extended recording
periods.
Rigid kinematic model assumptions about the chain between individual sensor sites
frequently diverge from physiological reality.
More critically, standard \acs{IMU} fusion algorithms use a magnetometer to correct yaw
drift.
In environments containing ferromagnetic structures --- steel reinforcement, hospital
beds, electric motors, or computer hardware --- local field distortion renders the
compass dysfunctional and introduces unbounded yaw error within minutes
\cite{madgwick2011}.

\subsection{Static clinical surface measurement}

Instruments such as the Idiag M360 (SpinalMouse) provide sagittal and frontal
curvature data with high accuracy (published precision \ang{1}--\ang{2};
\cite{spinalMouse}) and without radiation exposure.
They are well suited to supervised in-clinic assessments and therapeutic follow-up, but
by design they yield a single static snapshot.
They cannot record what happens to the lumbar spine during eight hours of warehouse work
or across a night's sleep.
